\section{Identification software must preserve what is and is not identified}
\label{exec18:r-identification}

Suppose a bank is considering an encouragement to make an additional repayment. Let $A$ denote randomized access to a reminder and $Y$ denote repayment by twelve months. The business question may ask for the effect of offering the reminder, whether a particular customer would repay only if reminded, or the range of possible repayment effects when an assumption is unavailable. Those are three different targets. A fitted conditional association, an interventional functional, a cross-world probability and a sharp identified set do not have interchangeable meanings. The proposed interface must name the target and information regime before it calls an R package.

For a supplied directed acyclic graph, \texttt{dagitty} can enumerate sufficient adjustment sets under its documented graphical criteria. That is a useful graph diagnostic; it is not a general algorithm for deriving every interventional functional under latent confounding. The existing \texttt{pcalg} and \texttt{bnlearn} comparison adds equivalence-class and causal-network procedures, but their returned graph or effect list has its own model class and assumptions. In particular, \texttt{pcalg}'s IDA manual assumes a faithful multivariate Gaussian law and a causal DAG without hidden variables, and estimates total effects over admissible DAG extensions. It is neither an arbitrary nonlinear effect region nor a confidence set for a bank decision. The \texttt{ivmodel} confidence-set interface answers another question again: its documented AR/CLR inversion returns interval components under its declared IV design, with CLR described as robust to weak instruments. The proposed wrapper must preserve these set components and admit infinite endpoints where the inversion yields them. A single label such as “causal effect” would erase these distinctions~\citep{exec13_other_dagitty,exec13_pcalg,exec13_bnlearn,v12_sw_ivmodel}.

The R package \texttt{causaleffect} supplies symbolic ID/IDC functionality and selected transportability and selection-recovery procedures. Its pinned R function accepts outcome, treatment, optional conditioning variables and a graph; it can return an expression or a derivation object. The selected implementation has an important adapter consequence: with its default \texttt{stop\_on\_nonid=TRUE}, nonidentification raises an error; with the option false, a failed identification can be represented by \texttt{id=FALSE} in the step result, an empty string, or \texttt{NULL}, depending on the other output options. A production wrapper must distinguish that documented algorithmic result from parse errors, unsupported graphs and runtime exceptions. It should also retain the graph class and functional, not just a Boolean or formatted string~\citep{exec18_rid_causaleffect}.

The pinned \texttt{dosearch} package asks whether an interventional target follows from multiple observational or experimental distributions and a graph, including selected transport, selection and missingness mechanisms. Its 2021 paper describes a sound search over do-calculus rules, with completeness tied to the specified problem class. Its Table 1 distinguishes complete ID/IDC cases, unknown completeness for general multiple-source identification, and an incomplete missing-data extension. The package documentation expressly says its missing-data search is not complete. More subtly, the pinned C++ source applies a positive time limit by returning early from the search; its initializer then decides identifiability from whether a target distribution was found and does not expose a separate timeout state in that result. Therefore, under that code path, a limit-truncated search can be represented like “not identifiable by this search.” An adapter must not promote that output to a proof of nonidentifiability. It must either disable the limit for the bounded admitted problem, instrument a distinct timeout/unknown state, or report the uncompleted search as unknown. This is a source-level review of a pinned revision, not an executed package trial~\citep{exec18_rid_dosearch,exec18_rid_dosearch_paper}.

Counterfactual identification is a further level. \texttt{cfid} implements ID*/IDC* for counterfactual conjunctions by constructing a parallel-worlds graph with shared latent variables and then applying recursive identification. Its function defaults to identifying from interventions, not from observations; the observational level must be requested or obtained through the documented alternatives. A true \texttt{id} field is not sufficient on its own: an inconsistent conjunction may be identifiable with probability zero, and a conditional query may be undefined when its conditioning event has zero probability. An adapter therefore retains the requested data level, identifying formula, returned \texttt{undefined} field and query itself; an inconsistency represented by a zero formula must retain that meaning. The package paper explains why shared latent variables connect the counterfactual worlds; treating them as independent simulated populations would change the query~\citep{exec18_rid_cfid,exec18_rid_cfid_paper}.

When point identification fails, bounds may still be useful, but their scope is not universal. \texttt{causaloptim} encodes certain graph-defined response-function problems as linear programs and permits added constraints; its documented graph attributes and admissible left/right structure delimit the class. \texttt{bpbounds} implements binary treatment and outcome bounds for an instrument with two or three categories, in a single trivariate table or a specified two-source bivariate design, with separate monotonicity outputs. Neither manual claims to deliver a sampling confidence region for every reported plug-in endpoint. The distinction matters: an identified set describes values compatible with a population law and assumptions; sampling uncertainty concerns how an estimator or set varies across repeated samples. A posterior interval from a separate Bayesian model is different again. The software contract should return these objects with different types and labels, including empty, unbounded, undefined and unknown states where relevant~\citep{exec18_rid_causaloptim,exec18_rid_bpbounds}.

\texttt{CausalQueries} is a useful adjacent tool rather than a substitute for a graphical identification test. Its documentation defines binary causal models, updates model parameters and supports queries over priors, parameters or posterior draws. Its query manual calls the default interval a “confidence interval” under the statistics argument, but names the interval endpoints as credible summaries and allows a prior or posterior distribution as the input. The proposed interface must attach the distribution and uncertainty meaning explicitly instead of inheriting that ambiguous label. Such a posterior answers a question conditional on the declared graph, response-type parameterization, restrictions, prior and update model. It need not equal a likelihood-identified set, and a narrow credible interval cannot certify that omitted graphs or cross-world dependence have been resolved. For the proposed library, the package is a reference for probabilistic model-query syntax and an optional R adapter candidate. The described package and its compiled Stan dependencies have not been installed, executed or cleared for redistribution in this review~\citep{exec18_rid_causalqueries}.

\subsection{An exact repayment example separates four outputs}

Let $Y_0,Y_1\in\{0,1\}$ be repayment under no reminder and reminder. Randomization identifies $\Pr(Y_0=1)=1/4$ and $\Pr(Y_1=1)=3/4$ in this deliberately stipulated population example. Write $q_{ij}=\Pr(Y_0=i,Y_1=j)$ and let $t=q_{11}$. The observed arm risks imply
\begin{align}
 q_{11}&=t, &q_{10}&=\frac14-t,\notag\\
 q_{01}&=\frac34-t, &q_{00}&=t,\qquad 0\leq t\leq\frac14. \label{exec18:r-fiber}
\end{align}
The intervention effect is point identified throughout this fiber:
\begin{equation}
 \operatorname{ATE}=\mathbb E(Y_1-Y_0)=\frac34-\frac14=\frac12.
 \label{exec18:r-ate}
\end{equation}
But the retrospective “necessity” fraction among would-be responders under the reminder is
\begin{equation}
 \operatorname{PN}=\Pr(Y_0=0\mid Y_1=1)
 =\frac{q_{01}}{q_{01}+q_{11}}=1-\frac{4t}{3}\in\left[\frac23,1\right].
 \label{exec18:r-pn}
\end{equation}
The average effect is a point, while the counterfactual fraction is only set identified. Assuming monotonicity $Y_1\geq Y_0$ removes the response type $q_{10}$, forces $t=1/4$ and narrows PN to $2/3$. That is an added scientific restriction, not information generated by a software call.

Every value in this interval is attainable: choose the corresponding $t$ and sample the four response types with probabilities $q_{ij}$, independently of randomized assignment. This supplies a constructive sharpness argument. The population risks above are stipulated constants; observing one success versus three in two samples of four would not establish those risks without sampling uncertainty.

The same distinction survives Bayesian updating. Consider the three compatible worlds $t\in\{0,1/8,1/4\}$. With four randomized units in each arm and the stated one-of-four versus three-of-four successes, the likelihood of the ordered outcomes is $(1/4)^2(3/4)^6=729/65536$ in every world. Thus the data do not change their relative weights. A uniform prior on these three worlds gives posterior mean PN $5/6$; prior weights $(1/6,1/3,1/2)$ give $7/9$. These are exact project calculations, not output from \texttt{CausalQueries} and not a defect in Bayesian computation: the difference records prior sensitivity along a likelihood-flat fiber. The aggregate binomial likelihood is $729/4096$ in every world, including the common combinatorial factor.

For an illustrative retrospective accounting rule, assign USD 1,000 to a responder whose repayment is attributed to the reminder and subtract USD 800 per responder. The compatible net amount per responder ranges from $1000(2/3)-800=-\$400/3$ to $1000(1)-800=\$200$. The two posterior means are USD $1000(5/6)-800=100/3$ and USD $1000(7/9)-800=-200/9$, respectively. These stipulated amounts are not observed economics or ROI, and attribution after the outcome must not be misread as prospective treatment selection on future repayment. The executable Fraction reference enumerates all three response-type distributions and checks every displayed likelihood and monetary value~\citep{exec18_rid_exact}.

\subsection{Conditional reuse and the boundary of a negative answer}

The lending case makes the graph restriction concrete. Suppose unrecorded risk $U$ affects both the offer $A$ and repayment $Y$. In a deliberately simple example, take $U\sim\operatorname{Bernoulli}(1/2)$ and $A=U$. Two structural models, $Y=A$ and $Y=U$, give the same observed $A=Y$ distribution, yet their average effects of intervening on $A$ are respectively one and zero. The graph allowing $U\to A$, $U\to Y$ and $A\to Y$ contains both models, because an allowed parent may have zero effect. No precision of fitting that observed distribution selects the intervention answer. This is an identification failure relative to that information and graph class. The earlier randomized example changes the information and identifies the average effect, while leaving cross-world attribution ambiguous.

For this use, the proposed adapter result retains three epistemic states: identified with a functional or point; not identified by a complete algorithm for its declared graph/data class; and unknown because the class is unsupported or computation ended before a certificate. Invalid input is a separate validation failure, with no scientific conclusion. Sharp identified sets, sampling confidence sets and Bayesian posterior summaries have separately named result types and assumptions. A zero-probability conditioning event is an undefined query, not a negative identification certificate.

The practical proposal is to retain \texttt{causaleffect} as a candidate symbolic R service and \texttt{cfid} as a counterfactual-query service, preserving their query objects and derivations. \texttt{dosearch} complements those services for multiple source distributions, but an adapter must record admitted search rules and termination before interpreting a negative result. Disabling a time limit does not make an incomplete rule system complete. Porting any of these algorithms requires an independently checked mathematical contract and differential fixtures; matching a formatted expression alone would miss a changed conditioning event or data level. For that reason direct source ports are deferred while optional process-bound R adapters are designed. This is a proposed interface boundary, not an executed integration or a conclusion about license compatibility~\citep{exec18_rid_causaleffect,exec18_rid_cfid,exec18_rid_dosearch}.

The restricted bounds engines offer a different reuse choice. \texttt{causaloptim} is a candidate optimization reference when its response-function class is appropriate; \texttt{bpbounds} supplies a smaller binary-IV reference with explicit monotonicity and two-source assumptions. The proposal rejects either as a universal sharp-bounds backend. It retains \texttt{CausalQueries} as an optional model-query path when priors and sensitivity are substantively justified, while rejecting a posterior interval as a replacement for identification. Dependency and license observations inform these choices but do not qualify redistribution, linked dependencies, numerical behavior or installed package execution~\citep{exec18_rid_causaloptim,exec18_rid_bpbounds,exec18_rid_causalqueries}.

These comparisons establish distinct candidate roles and failure semantics. Coverage of the wider R ecosystem remains incomplete; selected packages and their inspected versions do not exhaust the available methods. The exact examples establish the stated mathematical distinctions and illustrative accounting consequences, while implementation and broader source coverage remain future work.
